\Needspace{3.5in}
\begin{prob}

    Consider the following function which attempts to compute the maximum of a list of
    numbers.

    \begin{minted}[autogobble]{python}
        import math
        def maximum(numbers):
            n = len(numbers)
            if n == 0:
                return 0
            middle = math.floor(n / 2)
            return (
                maximum(numbers[:middle])
                +
                maximum(numbers[middle:])
            )
    \end{minted}

    Is this code guaranteed to work for all non-empty lists?

    You should adopt the convention that the maximum of an empty list is $-\infty$.

    \begin{choices}
        \choice Yes: it will return the maximum number in the list.
        \correctchoice No: it may recurse infinitely.
        \choice No: it may raise an exception.
        \choice No: it will run without error, but it might return the wrong answer.
    \end{choices}
\end{prob}
